% =============================================================================
%  MARCO TEÓRICO — Medición de Dispositivos Bajo Prueba (DUTs)
%  Calibración de Analizador de Redes Vectorial
%  Proyecto Final de Ingeniería Electrónica — UTN FRBA
%  Axel Nahum — 2026
% =============================================================================

\section{Marco Teórico: Medición de Parámetros S y Calibración del VNA}

El presente capítulo desarrolla los fundamentos teóricos de la medición de
parámetros S mediante un analizador de redes vectorial (VNA), los modelos de
error que afectan las mediciones reales, y los cinco métodos de calibración
implementados en el sistema NanoVNA-UTN-Toolkit para la caracterización de
dispositivos bajo prueba (DUTs).

\newpage
\Encabezado

% -----------------------------------------------------------------------------
\subsection{Parámetros S}
\label{sec:sparams}
% -----------------------------------------------------------------------------

Los parámetros de dispersión (S-parameters) son la representación matricial
más adecuada para caracterizar redes de dos puertos en alta frecuencia. A
diferencia de los parámetros Z, Y o H, los parámetros S se miden con los
puertos terminados en la impedancia de referencia $Z_0 = 50~\Omega$, condición
fácilmente realizable en la práctica.

Para una red de dos puertos, la matriz S se define por:

\begin{equation}
    \begin{pmatrix} b_1 \\ b_2 \end{pmatrix}
    =
    \begin{pmatrix} S_{11} & S_{12} \\ S_{21} & S_{22} \end{pmatrix}
    \begin{pmatrix} a_1 \\ a_2 \end{pmatrix},
    \label{eq:matriz_s}
\end{equation}

\noindent donde $a_i$ y $b_i$ son las ondas de tensión incidente y reflejada
normalizadas en el puerto $i$, respectivamente.

El significado físico de cada parámetro es:

\begin{itemize}
    \item $S_{11}$: coeficiente de reflexión en el puerto 1 con el puerto 2
          adaptado. Mide la energía que se refleja hacia la fuente.
    \item $S_{21}$: coeficiente de transmisión de puerto 1 a puerto 2.
          Mide la energía que atraviesa el dispositivo.
    \item $S_{22}$: coeficiente de reflexión en el puerto 2 (análogo a $S_{11}$).
    \item $S_{12}$: coeficiente de transmisión inverso (de puerto 2 a puerto 1).
\end{itemize}

En módulo, $|S_{11}|^2$ representa la fracción de potencia reflejada y
$|S_{21}|^2$ la fracción transmitida. Se expresan habitualmente en decibeles:

\begin{equation}
    S_{11}[\text{dB}] = 20\log_{10}|S_{11}|, \qquad
    S_{21}[\text{dB}] = 20\log_{10}|S_{21}|.
\end{equation}

Para un dispositivo pasivo de dos puertos sin pérdidas se cumple
$|S_{11}|^2 + |S_{21}|^2 = 1$ (conservación de la energía).

\newpage
\Encabezado

% -----------------------------------------------------------------------------
\subsection{El Analizador de Redes Vectorial (VNA)}
\label{sec:vna}
% -----------------------------------------------------------------------------

Un VNA mide los parámetros S de un dispositivo inyectando una señal de
prueba sinusoidal en un puerto y midiendo simultáneamente las ondas incidente
y reflejada/transmitida en todos los puertos. La medición es vectorial porque
captura tanto la magnitud como la fase de cada parámetro S.

La NanoVNA, empleada en este sistema, es un VNA de bajo costo de dos puertos
que opera en el rango de 50~kHz a 3~GHz (según modelo). Internamente realiza
una conversión de frecuencia hacia abajo (\emph{downconversion}) para procesar
la señal a una frecuencia intermedia inferior, donde la digitalización es más
sencilla.

\subsubsection{Errores sistemáticos de un VNA}

Las mediciones de un VNA real difieren de los valores verdaderos del DUT debido
a imperfecciones del instrumento, del cableado y de los conectores. Estos
errores son \emph{sistemáticos} (reproducibles y deterministas) en lugar de
aleatorios, por lo que pueden modelarse y corregirse mediante calibración.

Los principales errores sistemáticos son:

\begin{itemize}
    \item \textbf{Directividad} ($e_{00}$): señal que llega al receptor de
          reflexión sin haber interactuado con el DUT, debido a acoplamientos
          parásitos en el divisor de potencia interno del VNA.
    \item \textbf{Adaptación de fuente} ($e_{11}$, \emph{source match}):
          reflexiones producidas por la impedancia no perfectamente igual a
          $Z_0$ en el puerto de excitación, que se reinyectan en el DUT y
          producen múltiples rebotes.
    \item \textbf{Seguimiento de reflexión} ($e_{10}e_{01}$,
          \emph{reflection tracking}): variación en frecuencia de la ganancia
          del camino de la señal de prueba y del receptor de reflexión.
    \item \textbf{Seguimiento de transmisión} ($e_{10}e_{32}$,
          \emph{transmission tracking}): ídem para el camino de transmisión.
    \item \textbf{Adaptación de carga} ($e_{22}$, \emph{load match}):
          reflexiones en el puerto de recepción debidas a su impedancia no
          ideal.
\end{itemize}

\newpage
\Encabezado

% -----------------------------------------------------------------------------
\subsection{Modelo de Error del VNA}
\label{sec:modelo_error}
% -----------------------------------------------------------------------------

El modelo de error más utilizado para un VNA de dos puertos es el de
\emph{grafo de flujo de señal} (\emph{signal flow graph}). Para la medición
de reflexión (dirección directa), el modelo de 3 términos para S11 es
suficiente cuando solo se requiere corregir la reflexión.

\subsubsection{Modelo de 3 términos (1-puerto)}

El modelo de flujo de señal para la medición de $S_{11}$ incluye tres términos
de error~\cite{pozar2012, agilent2002}:

\begin{itemize}
    \item $e_{00}$: directividad.
    \item $e_{11}$: adaptación de fuente.
    \item $e_{10}e_{01}$: seguimiento de reflexión.
\end{itemize}

La relación entre el $S_{11}$ medido ($S_{11,m}$) y el $S_{11}$ real del DUT
($S_{11,c}$) queda:

\begin{equation}
    S_{11,m} = e_{00} + \frac{e_{10}e_{01}\cdot S_{11,c}}{1 - e_{11}\cdot S_{11,c}}.
    \label{eq:s11_medido}
\end{equation}

Definiendo $\Delta_e = e_{00}\cdot e_{11} - e_{10}e_{01}$, la corrección
inversa (fórmula de calibración SOL) es:

\begin{equation}
    \boxed{
    S_{11,c} = \frac{S_{11,m} - e_{00}}{S_{11,m}\cdot e_{11} - \Delta_e}.
    }
    \label{eq:correccion_s11}
\end{equation}

\subsubsection{Modelo de respuesta en transmisión (1 término)}

Para la corrección de $S_{21}$ mediante una normalización simple, el modelo
reduce a un único término de error de transmisión $e_{10}e_{32}$:

\begin{equation}
    S_{21,m} = e_{10}e_{32}\cdot S_{21,c},
    \label{eq:s21_medido}
\end{equation}

cuya corrección es trivialmente:

\begin{equation}
    \boxed{
    S_{21,c} = \frac{S_{21,m}}{e_{10}e_{32}}.
    }
    \label{eq:correccion_s21_simple}
\end{equation}

\subsubsection{Modelo de respuesta mejorada en transmisión (3 términos)}

El método de respuesta mejorada (\emph{Enhanced-Response}) extiende el modelo
anterior incorporando la adaptación de carga $e_{22}$ para corregir
parcialmente el efecto de los múltiples rebotes entre el DUT y el puerto de
recepción~\cite{rytting1987}:

\begin{equation}
    S_{21,c} = \frac{S_{21,m}}{e_{10}e_{32}}
               \cdot
               \frac{e_{10}e_{01}}{e_{11}\cdot S_{11,m} - \Delta_e}.
    \label{eq:correccion_s21_enhanced}
\end{equation}

Esta expresión incluye el denominador del modelo OSM, lo que permite compensar
la interacción entre la reflexión del DUT y la mismatch del puerto~2.

\newpage
\Encabezado

% -----------------------------------------------------------------------------
\subsection{Necesidad de Calibración}
\label{sec:necesidad_cal}
% -----------------------------------------------------------------------------

Sin calibración, los errores sistemáticos descriptos en la
Sección~\ref{sec:modelo_error} se superponen a la respuesta del DUT,
produciendo mediciones incorrectas. La magnitud de estos errores es
significativa: la directividad de un VNA de bajo costo como la NanoVNA
puede ser tan baja como 20~dB (en lugar de los 40--50~dB típicos de un VNA
de laboratorio), lo que introduce errores relativos de hasta el 10\,\% en
la magnitud de $S_{11}$ antes de cualquier calibración.

La calibración determina los coeficientes de error $\{e_{00}, e_{11},
e_{10}e_{01}, e_{10}e_{32}, e_{22}\}$ midiendo estándares de impedancia
\emph{conocidos} (circuito abierto, cortocircuito, carga adaptada, conexión
directa) y resolviendo el sistema de ecuaciones que relaciona los valores
medidos con los valores esperados.

Una vez determinados los coeficientes, la corrección se aplica en tiempo real
a cada barrido de frecuencia del DUT mediante las
fórmulas~\eqref{eq:correccion_s11}--\eqref{eq:correccion_s21_enhanced}.

La calibración debe rehacerse cuando cambia algún elemento del camino de
medición: se reemplaza el cable, se cambia la temperatura significativamente,
o se conecta/desconecta el cable del VNA.

\newpage
\Encabezado

% -----------------------------------------------------------------------------
\subsection{Estándares de Calibración}
\label{sec:estandares}
% -----------------------------------------------------------------------------

Los estándares de calibración son componentes de impedancia conocida que
permiten establecer la referencia de medición. En la banda de frecuencias
del sistema se utilizan los siguientes:

\begin{itemize}
    \item \textbf{Circuito abierto (Open)}: reflexión total con $\Gamma = +1$.
          Corresponde a un puerto sin conectar (o con un terminador de circuito
          abierto). El modelo ideal asume $S_{11} = +1\angle 0°$ a toda
          frecuencia.

    \item \textbf{Cortocircuito (Short)}: reflexión total con $\Gamma = -1$.
          Se logra cortocircuitando el conector. El modelo ideal asume
          $S_{11} = -1\angle 0°$ (o equivalentemente $1\angle 180°$) a toda
          frecuencia.

    \item \textbf{Carga adaptada (Match / Load)}: absorción total con
          $\Gamma = 0$. Se emplea una resistencia de $50~\Omega$ conectada
          al puerto. En el modelo ideal $S_{11} = 0$ y no hay reflexión.

    \item \textbf{Conexión directa (Thru)}: los dos puertos se conectan
          directamente (longitud cero). En el modelo ideal $S_{21} = 1\angle 0°$
          y $S_{11} = 0$.
\end{itemize}

En la práctica, los estándares reales se desvían del ideal a frecuencias
elevadas por inductancias y capacitancias parásitas del conector. Para el
rango de operación de la NanoVNA (hasta 3~GHz) con conectores SMA de buena
calidad, los modelos ideales resultan suficientemente precisos para la mayoría
de las aplicaciones.

\newpage
\Encabezado

% -----------------------------------------------------------------------------
\subsection{Métodos de Calibración Implementados}
\label{sec:metodos_cal}
% -----------------------------------------------------------------------------

El sistema NanoVNA-UTN-Toolkit implementa cinco métodos de calibración que
cubren distintos escenarios de medición, desde una simple normalización hasta
una calibración completa de reflexión y transmisión con compensación de
desadaptación de carga.

\subsubsection{OSM (Open -- Short -- Match)}
\label{sec:osm}

El método OSM es la calibración de 1-puerto estándar, conocida en la
literatura como SOL (Short -- Open -- Load)~\cite{pozar2012}.
Se emplean tres estándares y se corrige únicamente $S_{11}$.

\paragraph{Determinación de los coeficientes de error}

Midiendo los tres estándares ideales (Open: $\Gamma=+1$, Short: $\Gamma=-1$,
Match: $\Gamma=0$) se obtiene el sistema de tres ecuaciones que, resuelto en
forma cerrada, da~\cite{rytting1987}:

\begin{align}
    e_{00} &= S_{11,M},
    \label{eq:e00} \\[4pt]
    e_{11} &= \frac{S_{11,O} + S_{11,K} - 2\,e_{00}}{S_{11,O} - S_{11,K}},
    \label{eq:e11} \\[4pt]
    e_{10}e_{01} &= \frac{-2\,(S_{11,O} - e_{00})(S_{11,K} - e_{00})}
                         {S_{11,O} - S_{11,K}},
    \label{eq:e10e01}
\end{align}

\noindent donde $S_{11,O}$, $S_{11,K}$ y $S_{11,M}$ son los $S_{11}$ medidos
del Open, Short (Kurzschluss) y Match, respectivamente.

\paragraph{Corrección en tiempo real}

La corrección se aplica mediante~\eqref{eq:correccion_s11}:

\begin{equation}
    S_{11,c} = \frac{S_{11,m} - e_{00}}{S_{11,m}\cdot e_{11} - \Delta_e},
    \qquad \Delta_e = e_{00}\cdot e_{11} - e_{10}e_{01}.
    \nonumber
\end{equation}

\paragraph{Resumen}

\begin{itemize}
    \item Estándares: Open, Short, Match.
    \item Parámetro corregido: $S_{11}$.
    \item Coeficientes de error: $e_{00}$, $e_{11}$, $e_{10}e_{01}$.
    \item Aplicación típica: medición de impedancias, antenas, filtros (reflexión).
\end{itemize}

% - - - - - - - - - - - - - - - - - - - - - - - - - - - - - - - - - - - - - -
\subsubsection{Normalización Open/Short (Open/Short Normalization)}
\label{sec:os_norm}

Este método es una calibración de reflexión simplificada de 1 término.
Se realiza \emph{una única medición} de estándar (Open o Short) y se utiliza
directamente como coeficiente de seguimiento de reflexión.

\paragraph{Identificación automática del estándar}

El sistema detecta automáticamente si el estándar medido es Open o Short
evaluando el signo de la parte real media del $S_{11}$ medido:

\begin{equation}
    \text{estándar} =
    \begin{cases}
        \text{Short} & \text{si } \overline{\operatorname{Re}(S_{11})} < 0, \\
        \text{Open}  & \text{en otro caso.}
    \end{cases}
\end{equation}

Si se detecta un Short, el signo del $S_{11}$ almacenado se invierte para que
el coeficiente quede siempre en fase positiva.

\paragraph{Corrección}

\begin{equation}
    S_{11,c} = \frac{S_{11,m}}{e_{10}e_{01}},
    \label{eq:os_norm}
\end{equation}

seguida de un desenvuelto de fase (\emph{phase unwrapping}) para garantizar
continuidad de la respuesta en frecuencia.

\paragraph{Limitaciones}

Al no corregir directividad ni adaptación de fuente, este método solo elimina
la variación en frecuencia del camino de medición. Es adecuado para
dispositivos con reflexión moderada a alta ($|S_{11}| > -20$~dB), donde los
errores de directividad son pequeños relativos a la señal medida.

\paragraph{Resumen}

\begin{itemize}
    \item Estándares: Open \emph{o} Short (1 medición).
    \item Parámetro corregido: $S_{11}$ (normalización de amplitud y fase).
    \item Coeficientes de error: $e_{10}e_{01}$ únicamente.
    \item Aplicación típica: inspección rápida de reflexión, sin requerir carga adaptada.
\end{itemize}

% - - - - - - - - - - - - - - - - - - - - - - - - - - - - - - - - - - - - - -
\subsubsection{Normalización Thru (Thru Normalization)}
\label{sec:thru_norm}

Este método corrige únicamente $S_{21}$ mediante una sola medición de
transmisión con los puertos conectados directamente (Thru).

\paragraph{Calibración y corrección}

El $S_{21}$ medido del Thru se almacena como el coeficiente de seguimiento de
transmisión:

\begin{equation}
    e_{10}e_{32} = S_{21,\text{Thru}}.
\end{equation}

Durante la medición del DUT, la corrección es:

\begin{equation}
    S_{21,c} = \frac{S_{21,m}}{e_{10}e_{32}}.
    \label{eq:thru_norm}
\end{equation}

El parámetro $S_{11}$ no se corrige en este método.

\paragraph{Limitaciones}

La normalización Thru solo elimina la pérdida y el desfase del camino de
transmisión entre los puertos. No compensa la desadaptación de ninguno de los
puertos ni los múltiples rebotes entre el DUT y el sistema.

\paragraph{Resumen}

\begin{itemize}
    \item Estándares: Thru (1 medición de 2 puertos).
    \item Parámetro corregido: $S_{21}$.
    \item Coeficientes de error: $e_{10}e_{32}$ únicamente.
    \item Aplicación típica: medición de ganancia/pérdida de inserción.
\end{itemize}

% - - - - - - - - - - - - - - - - - - - - - - - - - - - - - - - - - - - - - -
\subsubsection{1-Port+N}
\label{sec:1portn}

El método 1-Port+N combina la calibración OSM completa para $S_{11}$ con la
normalización Thru para $S_{21}$, permitiendo corregir ambos parámetros en
un único flujo de calibración de 4 estándares.

\paragraph{Secuencia de calibración}

\begin{enumerate}
    \item Open → Short → Match: se calculan $e_{00}$, $e_{11}$, $e_{10}e_{01}$
          mediante~\eqref{eq:e00}--\eqref{eq:e10e01}.
    \item Thru: se almacena $e_{10}e_{32} = S_{21,\text{Thru}}$.
\end{enumerate}

\paragraph{Correcciones aplicadas}

\begin{align}
    S_{11,c} &= \frac{S_{11,m} - e_{00}}{S_{11,m}\cdot e_{11} - \Delta_e},
    \label{eq:1portn_s11} \\[4pt]
    S_{21,c} &= \frac{S_{21,m}}{e_{10}e_{32}}.
    \label{eq:1portn_s21}
\end{align}

\paragraph{Resumen}

\begin{itemize}
    \item Estándares: Open, Short, Match, Thru (4 mediciones).
    \item Parámetros corregidos: $S_{11}$ (3-términos) y $S_{21}$ (1-término).
    \item Aplicación típica: caracterización completa de dispositivos de 2 puertos
          donde la precisión en $S_{21}$ no es crítica.
\end{itemize}

% - - - - - - - - - - - - - - - - - - - - - - - - - - - - - - - - - - - - - -
\subsubsection{Respuesta Mejorada (Enhanced-Response)}
\label{sec:enhanced}

El método de respuesta mejorada extiende el 1-Port+N incorporando la
corrección de la adaptación de carga $e_{22}$ en la medición de $S_{21}$.
Utiliza los mismos 4 estándares que el 1-Port+N, pero extrae información
adicional de la medición del Thru.

\paragraph{Cálculo de los coeficientes adicionales}

A partir de la medición del Thru (con $S_{11}$ y $S_{21}$ conocidos) y los
coeficientes OSM ya calculados se obtienen:

\begin{align}
    e_{22} &= \frac{S_{11,\text{Thru}} - e_{00}}
                   {S_{11,\text{Thru}}\cdot e_{11} - \Delta_e},
    \label{eq:e22} \\[4pt]
    e_{10}e_{32} &= S_{21,\text{Thru}}\,(1 - e_{11}\cdot e_{22}).
    \label{eq:e10e32}
\end{align}

\paragraph{Corrección mejorada de $S_{21}$}

\begin{equation}
    S_{21,c} = \frac{S_{21,m}}{e_{10}e_{32}}
               \cdot
               \frac{e_{10}e_{01}}{e_{11}\cdot S_{11,m} - \Delta_e}.
    \label{eq:enhanced_s21}
\end{equation}

El segundo factor en~\eqref{eq:enhanced_s21} compensa la interacción entre
la reflexión de entrada del DUT y la desadaptación del puerto de recepción,
reduciendo el error de medición en dispositivos con alta reflexión.

\paragraph{Resumen}

\begin{itemize}
    \item Estándares: Open, Short, Match, Thru (4 mediciones, mismos que 1-Port+N).
    \item Parámetros corregidos: $S_{11}$ (3-términos) y $S_{21}$ (3-términos).
    \item Coeficientes: $e_{00}$, $e_{11}$, $e_{10}e_{01}$, $e_{22}$, $e_{10}e_{32}$.
    \item Aplicación típica: caracterización de dispositivos de 2 puertos con
          mayor precisión en $S_{21}$.
\end{itemize}

\newpage
\Encabezado

% -----------------------------------------------------------------------------
\subsection{Kits de Calibración}
\label{sec:kits}
% -----------------------------------------------------------------------------

El sistema permite guardar una calibración completa como un \emph{kit},
identificado por un nombre y asociado a un método y fecha de medición.
Un kit almacena:

\begin{itemize}
    \item Las mediciones crudas de los estándares (archivos Touchstone .s1p,
          .s2p).
    \item Los coeficientes de error calculados, también en formato Touchstone.
    \item Metadatos: nombre del kit, método de calibración, fecha y hora.
\end{itemize}

Los kits permiten reutilizar una calibración sin necesidad de volver a medir
los estándares, siempre que el camino de medición (cables, conectores, VNA)
no haya cambiado desde la última calibración. Esta funcionalidad es
especialmente útil cuando se realizan series largas de mediciones del mismo
tipo de DUT.

\newpage
\Encabezado

% -----------------------------------------------------------------------------
\subsection{Resumen Comparativo de los Métodos de Calibración}
\label{sec:resumen_cal}
% -----------------------------------------------------------------------------

La Tabla~\ref{tab:comparacion_cal} resume las características de los cinco
métodos implementados.

\begin{table}[h!]
\centering
\caption{Comparación de los métodos de calibración del sistema NanoVNA-UTN-Toolkit.}
\label{tab:comparacion_cal}
\begin{tabular}{@{}lccccc@{}}
\toprule
\textbf{Método} &
\textbf{Estándares} &
\textbf{Términos} &
$S_{11}$ &
$S_{21}$ &
\textbf{Uso típico} \\
\midrule
OSM              & O, K, M        & 3 & Completo   & ---         & Reflexión \\
Open/Short Norm. & O \emph{o} K   & 1 & Normaliz.  & ---         & Reflexión rápida \\
Thru Norm.       & Thru           & 1 & ---         & Normaliz.   & Ganancia/pérdida \\
1-Port+N         & O, K, M, Thru  & 4 & Completo   & Normaliz.   & 2-puertos general \\
Enhanced-Resp.   & O, K, M, Thru  & 5 & Completo   & Mejorado    & 2-puertos preciso \\
\bottomrule
\end{tabular}
\end{table}

\noindent Referencias: O = Open, K = Short (\emph{Kurzschluss}), M = Match.

% -----------------------------------------------------------------------------
%  REFERENCIAS
% -----------------------------------------------------------------------------
%
% \bibitem{pozar2012}
% D.~M.~Pozar,
% \emph{Microwave Engineering}, 4th ed.
% John Wiley \& Sons, Hoboken, NJ, 2012.
%
% \bibitem{rytting1987}
% D.~Rytting,
% ``ARFTG short course notes: Network analyzer error models and calibration methods,''
% Agilent Technologies (HP), 1987 (revisado 2001).
%
% \bibitem{agilent2002}
% Agilent Technologies,
% \emph{Applying Error Correction to Network Analyzer Measurements},
% Application Note AN~1287-3, 2002.
%
